\section{De las islas de características manuales a una arquitectura unificada}

La sección anterior explicó que esta clase busca responder «por qué existe un curso sobre Transformer». Esta sección da la respuesta histórica: hace más de diez años, cada área no solo usaba datos distintos, sino también características, vocabularios y pilas de software distintos; el aprendizaje profundo unificó primero el paradigma de entrenamiento, y Transformer unificó después la arquitectura troncal. La unificación redujo el costo de aprender de un área a otra y convirtió el cómputo, los datos y la escala en variables comunes.

\subsection{2011: la ingeniería de características ensamblada en un pipeline frágil}

Karpathy parte de su propia llegada a la visión por computadora. El problema no era que los investigadores «no fueran inteligentes», sino que el sistema delegaba el diseño de la representación en una gran cantidad de descriptor manuales: cada artículo aportaba una característica local, y los practicantes las extraían todas, las concatenaban y entrenaban un clasificador encima. Este flujo era difícil de optimizar de forma conjunta y también hacía difícil determinar si un error provenía de las características, de los datos o del clasificador.

\lecturefigure{slide-11-historical-context.jpg}{Contexto histórico: una mirada a 2023 desde el flujo de trabajo de visión de 2011.}{Video oficial de Stanford Online, 00:10:39--00:11:17.}

\teachervoice{Se compara con humor con Bilbo, el personaje de «El Señor de los Anillos» que rememora el pasado, y pregunta a los estudiantes si se dan cuenta de lo afortunados que son al entrar a la IA en 2023. Detrás del humor hay un juicio: la cadena de herramientas unificada que hoy damos por sentada no existió de forma natural, sino que es el resultado gradual del aprendizaje de representaciones, el entrenamiento a gran escala y las arquitecturas compartidas.}

\lecturefigure{slide-12-scene-2011.jpg}{«Scene: it is 2011»: la entrada a la época de los pipelines de visión manuales.}{Video oficial de Stanford Online, 00:11:17--00:11:30.}

Una \term{handcrafted feature} (característica manual) es un estadístico de la imagen definido de antemano por el investigador, como gradientes locales, texturas, histogramas de color o plantillas geométricas; un \term{descriptor} es la regla que codifica una región local como un vector. Contienen un inductive bias fuerte, pero no permiten que todas las etapas se ajusten en conjunto en torno a la pérdida final, como sí lo hace una red neuronal de extremo a extremo.

\lecturefigure{slide-13-old-computer-vision-paper.jpg}{Varias páginas de características y detalles de ingeniería en un artículo de visión de la época anterior.}{Video oficial de Stanford Online, 00:11:30--00:12:11.}

\begin{knowledgebox}{Lectura de la figura: no confundir la densidad de términos con «un modelo más profundo»}
La densidad de la página proviene sobre todo del feature zoo: cada descriptor extrae por su cuenta bordes, texturas, colores o relaciones espaciales, y luego los vectores se envían a una SVM. Una \term{SVM} (Support Vector Machine, máquina de vectores de soporte) es un modelo supervisado que busca una frontera de clasificación en el espacio de características. Puede ser muy potente, pero la representación de las etapas previas suele estar fija, y el pipeline completo no es un sistema diferenciable optimizado de extremo a extremo por un mismo objetivo.
\end{knowledgebox}

La crítica aguda de Karpathy es esta: el código se reunía de todas partes, y el sistema no solo era complejo, sino que producía predicciones absurdas que hoy parecerían bugs. En aquel entonces, los practicantes a veces las consideraban fallas ocasionales inevitables. Este contraste refleja un cambio en la cultura de ingeniería: cuando los modelos, los datos y las rutas de entrenamiento están más unificados, es más probable que los errores se localicen, se reproduzcan y se corrijan de manera sistemática.

\lecturefigure{slide-14-old-lstm-paper.jpg}{Artículos antiguos de otra área: los modelos de secuencias tenían su propia terminología, sus propios diagramas y su propia tradición de implementación.}{Video oficial de Stanford Online, 00:12:23--00:12:56.}

Los artículos sobre LSTM y los de visión no solo trataban objetos distintos, sino que también descomponían los problemas de manera distinta. Cuando un investigador de visión veía hidden state, gates y time unrolling, tenía que aprender de nuevo todo un lenguaje; un investigador de secuencias enfrentaba el mismo obstáculo al leer sobre descriptor de visión y pipelines geométricos. Por ello, el costo de transferir conocimiento entre áreas no provenía solo de la dificultad matemática, sino también de la falta de una interfaz conceptual común.

\lecturefigure{slide-15-old-language-model-paper.jpg}{Artículo antiguo sobre modelos de lenguaje: otro conjunto de tareas, notación y tradición de evaluación.}{Video oficial de Stanford Online, 00:12:23--00:12:56.}

El objetivo central de un modelo de lenguaje puede ser muy uniforme: predecir el siguiente token a partir del contexto. Sin embargo, en los primeros artículos, los n-gramas, el suavizado, la morfología, la sintaxis y las características específicas de cada tarea solían formar pilas técnicas independientes. Uno de los avances del aprendizaje profundo fue hacer que el embedding y el predictor se aprendieran juntos a partir de los datos, de modo que «cómo se representa el contexto» dejara de depender por completo de una enumeración manual.

\lecturefigure{slide-16-old-neural-architecture.jpg}{Diagrama de una arquitectura neuronal antigua: las redes de cada área aún mostraban lenguajes estructurales muy distintos.}{Video oficial de Stanford Online, 00:12:23--00:12:56.}

El punto didáctico de esta figura no es identificar cada recuadro, sino observar las rutas de la información. Los sistemas tempranos solían codificar los supuestos de la tarea directamente en la topology: conexiones bidimensionales locales para imágenes, recurrencia temporal para texto y etapas de entrada especializadas para voz. Ese bias puede mejorar la eficiencia de muestreo, pero obliga a rediseñar el esqueleto ante una nueva modalidad, un nuevo sensor o relaciones de largo alcance.

\lecturefigure{slide-17-ai-areas-different.jpg}{Antes de 2012: la visión, el lenguaje, la voz, la traducción y el RL parecían disciplinas distintas.}{Video oficial de Stanford Online, 00:12:56--00:13:17.}

\begin{warningbox}{Unificar no equivale a borrar el conocimiento de cada área}
Un lenguaje común de redes neuronales reduce la barrera de lectura, pero no implica que todas las áreas requieran el mismo procesamiento de datos. La vecindad bidimensional de las imágenes, la estructura tiempo-frecuencia de la voz, las condiciones de acción del RL y las restricciones geométricas de las proteínas todavía deben incorporarse en la tokenización, el position/type embedding, la mask, el objective o un specialized module. Lo que se unifica es un tronco componible, no el mundo real en sí.
\end{warningbox}

\subsection{2012: el aprendizaje a gran escala unifica el paradigma de entrenamiento}

El problema anterior era que cada área definía sus representaciones a mano. La época de ImageNet mostró otra ruta: grandes conjuntos de datos, modelos grandes y hardware acelerado podían permitir que la red neuronal aprendiera por sí misma representaciones jerárquicas. Así, el foco de la investigación pasó de «inventar otro descriptor» a «cómo obtener datos, cómputo y una optimización estable», y esta recipe se replicó después en voz, traducción y RL.

\lecturefigure{slide-18-imagenet-2012.jpg}{ImageNet 2012: la escala, los datos y el entrenamiento en GPU cambiaron el foco de la investigación en visión.}{Video oficial de Stanford Online, 00:12:56--00:13:33.}

\begin{knowledgebox}{Lectura de la figura: lo que realmente cambió fue el ciclo de optimización}
Los píxeles de entrada pasan por una red de varias capas con parámetros compartidos para producir una predicción, y la pérdida de clasificación ajusta todas las capas mediante backpropagation. La representación deja de ser una interfaz manual congelada y se aprende junto con el objetivo de la tarea. Lo decisivo aquí no es solo la convolución en sí, sino que «modelo diferenciable + datos a gran escala + cómputo en paralelo + pérdida de extremo a extremo» forman un ciclo replicable.
\end{knowledgebox}

\lecturefigure{slide-19-areas-read-more-similar.jpg}{Después del aprendizaje profundo: los artículos de distintas áreas empiezan a compartir parámetros, optimizadores y lenguaje de entrenamiento.}{Video oficial de Stanford Online, 00:13:33--00:14:10.}

Karpathy subraya que de pronto apareció un vocabulario común reconocible en artículos de distintas áreas: parameters, optimizer, loss, neural network. Un marco de representación unificado produce reutilización del conocimiento: los investigadores de visión pueden entender los problemas de optimización de la traducción automática, y los investigadores de NLP pueden tomar las conexiones residuales y el entrenamiento a gran escala. Una cadena de herramientas común adelgaza las fronteras entre áreas y crea las condiciones para la unificación arquitectónica posterior a 2017.

\subsection{2017: incluso la arquitectura troncal empieza a converger}

El aprendizaje profundo unificó primero «cómo entrenar», y Transformer unificó después «qué tipo de tronco entrenar». La aproximación de primer orden de Karpathy es esta: se hace copy-paste de la misma arquitectura a distintas áreas y lo que cambia principalmente es el chunking de los datos y la forma de entrada. La afirmación es exagerada a propósito, pero señala con precisión que la attention permite al modelo aprender relaciones dinámicas sobre un conjunto de elementos.

\lecturefigure{slide-20-transformer-2017.jpg}{Transformer en 2017: un artículo de traducción automática propone un paquete arquitectónico reutilizable.}{Video oficial de Stanford Online, 00:13:57--00:14:21.}

\begin{knowledgebox}{Lectura de la figura: la figura del artículo original no es solo attention}
Tanto el encoder de la izquierda como el decoder de la derecha contienen embedding, position, multi-head attention, feed-forward, residual add y normalization; el decoder requiere además causal mask y cross-attention. El título destaca la attention porque sustituye a la recurrence como principal vía de comunicación; que el modelo pueda entrenarse y escalarse depende de todo el package.
\end{knowledgebox}

\lecturefigure{slide-21-transformer-across-fields.jpg}{La misma figura de Transformer aparece una y otra vez en visión, lenguaje, voz, RL y aprendizaje sobre grafos.}{Video oficial de Stanford Online, 00:14:17--00:14:40.}

Esta diapositiva debe leerse como «capa de interfaz — capa compartida — capa de tarea»: la capa de interfaz asigna a vectores un patch, una word/subword, un spectrogram slice, un state/action/reward o un graph element; la capa compartida actualiza los vectores con attention y MLP; la capa de tarea interpreta la salida como clasificación, generación, control o predicción de estructura. La reutilización entre áreas ocurre en la parte intermedia; la entrada y el objetivo siguen determinando qué aprende el modelo.

\lecturefigure{slide-22-brain-unified-compute.jpg}{Analogía con la homogeneidad de la corteza: la atractiva conjetura de un algoritmo de aprendizaje unificado.}{Video oficial de Stanford Online, 00:14:42--00:15:09.}

\teachervoice{Karpathy dice que la corteza visual y la auditiva son similares a escala macroscópica, y que el hecho de que Transformer solo cambie algunos «hiperparámetros» entre áreas tal vez sugiera un algoritmo de aprendizaje unificado y potente. Es una analogía sugerente, no evidencia neurocientífica; el éxito de la transferencia arquitectónica puede deberse a una estructura de cómputo universal, pero también a una solución de ingeniería favorecida en conjunto por los datos y el hardware modernos.}

\subsection{Resumen de la sección}

La IA de 2011 estaba fragmentada por características manuales, vocabularios de cada área y topology especializadas; el aprendizaje profundo a gran escala posterior a 2012 unificó el lenguaje de la optimización de extremo a extremo; y en 2017 Transformer llevó el tronco reutilizable al nivel multimodal. La siguiente sección muestra que esta arquitectura no «cayó del cielo»: evolucionó a lo largo de una cadena de problemas que pasa por los modelos de lenguaje, el bottleneck de seq2seq y la soft alignment.
